% =====================================================================
\section{Spatial structure in neural operators and the U-Net}
\label{sec:spatial-structure}
% =====================================================================

Ionized bubbles span a wide range of sizes, while their boundaries can be sharp on the grid scale. This motivates three useful architectural choices: local interactions, a hierarchy of scales, and explicit modulation by the field. Existing models address these in different ways. This section reviews the U-Net and U-FNO before introducing the particular combination used by the rhizome.

% ---------------------------------------------------------------------
\subsection{Inductive biases: locality, scale and stationarity}
\label{sec:inductive-biases}

The spatial interaction in a neural-operator layer can often be written as a kernel integral, with pointwise maps around it. Restrictions on this interaction determine some of the architecture's inductive biases. Three properties are useful for comparison.

\paragraph{Receptive field.} A convolutional neural network (CNN) uses kernels with compact support of a few grid cells, $\kappa(r-r')=0$ for $\abs{r-r'}$ larger than the stencil. A single layer is therefore strictly local; the receptive field grows only linearly with depth. The FNO is the opposite extreme: every Fourier layer couples every point to every other point.

\paragraph{Scale.} A plain CNN or FNO processes the field at one resolution. Neither has a built-in notion of coarse and fine scales; scale-dependent behaviour has to be learned implicitly. Multi-scale architectures instead process the field at several resolutions and combine the results, so that large-scale context and small-scale detail are handled by separate parts of the network.

\paragraph{Stationarity and state dependence.} Convolutional filters and standard Fourier multipliers are fixed after training. Their surrounding nonlinearities still make the complete network's response depend on the input. Attention \citep{vaswani2017attention} and graph kernel networks \citep{li2020graph} put feature dependence directly into the interaction weights. Dense pairwise evaluation costs $\mathcal{O}(N^2)$ at fixed width, but this is not a lower bound for state dependence: sparse graph interactions and factorized or linear attention reduce that cost \citep{katharopoulos2020linear}. The rhizome uses a convolution with separable endpoint gates.

\Cref{tab:inductive-biases} summarizes these properties for the architectures discussed in this report.

\begin{table*}[!tp]
  \centering
  \caption{Spatial interactions in the architectures discussed here. ``Fixed'' refers to the learned linear filters, not to the full nonlinear response. Costs hold at fixed width; pooling and boundaries can break exact translation equivariance.}
  \label{tab:inductive-biases}
  \begin{tabular}{lcccc}
    \toprule
    Architecture & Receptive field & Base filter & Scale hierarchy & Explicit gates \\
    \midrule
    CNN          & local, grows with depth   & fixed & no  & no \\
    FNO          & global spectral path     & fixed & no  & no \\
    U-Net        & local $\to$ wide via pooling & fixed & yes & no \\
    U-FNO        & global + multi-scale local & fixed & yes & no \\
    Dense attention & global                  & adaptive & no & $\mathcal{O}(N^2)$ \\
    Rhizome      & global                    & fixed + gates & no & $\mathcal{O}(N\log N)$ \\
    \bottomrule
  \end{tabular}
\end{table*}

% ---------------------------------------------------------------------
\subsection{The U-Net}
\label{sec:unet}

The U-Net \citep{ronneberger2015unet} was introduced for biomedical image segmentation, a task that, like predicting an ionization field, requires sharp, pixel-accurate boundaries informed by large-scale context. It consists of three parts (\cref{fig:unet}).

\paragraph{Contracting path.} The encoder applies a stack of small convolutions followed by a nonlinearity, then halves the resolution by pooling, and repeats. Denoting the feature map at level $s$ by $e_s$,
\begin{equation}
  \label{eq:unet-encoder}
  e_0 = C_0(a), \qquad e_{s} = C_s\big(\mathrm{Down}(e_{s-1})\big), \quad s = 1,\dots,S,
\end{equation}
where each $C_s$ is a block of $3\times3$ convolutions with nonlinearities. The number of channels is typically doubled at each level, so the network trades spatial resolution for feature richness. Because every level operates on a grid that is coarser by a factor of two, a $3\times3$ convolution at level $s$ spans $3\cdot2^{s}$ cells of the original grid. The receptive field therefore grows exponentially with the number of levels instead of linearly with depth.

\paragraph{Bottleneck.} At the coarsest level, $e_S$, the field is represented on a grid of $n/2^{S}$ points per dimension, where $n$ is the original side length with many channels. Here a few convolutions see a large fraction of the domain and capture large-scale context.

\paragraph{Expanding path and skip connections.} The decoder upsamples step by step and, at each level, concatenates the upsampled features with the encoder features of the same resolution:
\begin{equation}
  \label{eq:unet-decoder}
  d_S = e_S, \qquad d_s = C'_s\big(\big[\,\mathrm{Up}(d_{s+1}),\ e_s\,\big]\big), \quad s = S-1,\dots,0,
\end{equation}
and a final $1\times1$ convolution maps $d_0$ to the output $u$. The skip connections $e_s\to d_s$ are the key ingredient. Pooling discards the exact position of features, and upsampling alone cannot restore it. The skip connections pass the full-resolution encoder features directly to the decoder, so the network can combine \emph{what} is present (from the coarse levels) with \emph{where} it is (from the fine levels).

\begin{figure}[!htbp]
  \centering
  \figplaceholder[\linewidth]{unet_architecture.pdf}{U-Net: contracting path,
  bottleneck, expanding path with skip connections.}
  \caption{U-Net architecture with $S=3$ downsampling steps. The contracting path (left) halves the resolution $N$ and doubles the channel count $c$ at each level; the expanding path (right) reverses this. Dashed skip connections concatenate encoder features with the upsampled decoder features of the same resolution, which restores the spatial detail lost by pooling.}
  \label{fig:unet}
\end{figure}

\paragraph{Why it captures sharp structure.} The U-Net provides an explicit hierarchy of local receptive fields. Its finest levels operate on the full grid with local kernels, so edges can be represented directly in physical space instead of through a truncated set of global modes. Its intermediate levels give it receptive fields at every scale between a few cells and the full box, which matches the broad range of bubble sizes. Finally, skip connections let high-frequency detail bypass the bottleneck entirely. U-Nets are therefore strong baselines for field-to-field emulation whenever sharp, localized features matter.

\paragraph{Limitations.} The convolution kernels are defined in grid cells, so their physical support changes with resolution. A fully convolutional U-Net can often run on a different grid size, with padding where needed, but this does not make it a consistent discretization of the same physical operator. Pooling also introduces alignment dependence, so exact equivariance to arbitrary one-cell shifts is generally lost. Large-scale context passes through coarse features; whether this is sufficient depends on the task. The linear filters are shared spatially, while the nonlinear response remains input-dependent.

% ---------------------------------------------------------------------
\subsection{Hybrid architectures: U-FNO and related}
\label{sec:ufno}

The complementary strengths of the FNO (global, resolution-independent) and the U-Net (local, multi-scale) suggest combining them. The U-FNO of \citet{wen2022ufno} does this inside the Fourier layer. In addition to the spectral path $\F^{-1}(R_\phi\cdot\F v)$ and the pointwise path $W v$ of \cref{fig:fno-architecture}, each \emph{U-Fourier layer} adds a small U-Net acting on $v$:
\begin{equation}
  \label{eq:ufno-layer}
  v_{\ell+1}(x) = \sigma\Big( W v_\ell(x) + \F^{-1}\big(R_\phi\cdot\F v_\ell\big)(x) + \mathrm{UNet}_\ell[v_\ell](x) \Big),
\end{equation}
where $\mathrm{UNet}_\ell$ denotes the U-Net path. The spectral path supplies global, low-$k$ structure, while the U-Net path restores the high-$k$, localized detail that mode truncation removes. In the original application to multiphase flow, where saturation fronts are sharp, this substantially improved accuracy over the plain FNO \citep{wen2022ufno}.

The improvement comes at a cost that matters when comparing architectures. The U-Net path adds a large number of parameters, and a U-FNO is typically much larger than an FNO of the same width and depth. In our own comparison (\cref{sec:results}), the U-FNO reference model has roughly two orders of magnitude more parameters than the matched models. Better accuracy can therefore not be attributed to the architecture alone. We treat that implementation as a dense reference baseline. Parameter matching alone is not enough to isolate architecture: inputs, losses, optimization and compute also need to be controlled. The convolutional path introduces cell-scale dependence into the U-FNO, even though its Fourier multipliers can be reused across grids.

Other hybrids follow the same pattern of adding a local or multi-scale branch next to a global spectral one. They differ in which branch is used and where it is inserted, but they combine spatial branches with fixed learned filters and nonlinear maps, as in \cref{eq:ufno-layer}.

% ---------------------------------------------------------------------
\subsection{Open problem: structure that adapts to the field}
\label{sec:adaptive-structure}

A useful question is whether an explicit, inexpensive modulation of the spatial kernel improves field prediction beyond what a stack of fixed filters and nonlinearities already learns. This is an architectural question, not a claim that standard nonlinear networks cannot adapt to their inputs.

Radiative transfer motivates such modulation: absorption depends on the intervening material, so a physical source--receiver response need not be stationary. The present targets, however, come from the semi-numerical prescriptions of \cmfast{}, not a resolved transport calculation. Neither attention based only on endpoint features nor the rhizome's endpoint gates automatically encode optical depth along a ray. Intermediate information can enter through repeated contextual updates, but a physical transport interpretation requires separate evidence.

An operator suited to this problem should combine four properties: (i) global reach, so that interactions across the whole box are possible; (ii) state-dependent coupling, so that the effective connections change as the field evolves; (iii) a cost close to $\mathcal{O}(N\log N)$, so that three-dimensional grids remain tractable; and (iv) a parameterization that is independent of the grid, so that the discretization invariance of the FNO is kept. The rhizome operator, introduced in \cref{sec:rhizome}, is designed around these four requirements.
